\chapter{Resumo}

O ENEM é a principal porta de entrada para o ensino superior no Brasil, porém os resultados mais recentes revelam queda persistente no desempenho em Matemática. Avaliações internacionais como o PISA 2022 confirmam esse cenário crítico: apenas 27\% dos estudantes brasileiros atingiram o nível básico em matemática \cite{brasil2026pisa}. Pesquisas apontam que essas dificuldades decorrem de estratégias cognitivas inadequadas, baixa compreensão dos enunciados e desengajamento. Apesar da ênfase da BNCC no Pensamento Computacional (PC) como habilidade fundamental para a resolução de problemas, poucas ferramentas educacionais integram efetivamente o PC a estratégias pedagógicas engajadoras. Este trabalho apresenta o \textbf{Cognitivo}, um software gamificado fundamentado nos quatro pilares do PC (decomposição, abstração, reconhecimento de padrões e algoritmos) projetado para aumentar o engajamento de estudantes em preparação para a prova de matemática do ENEM. Espera-se que a integração entre gamificação e PC estimule o engajamento, o raciocínio lógico e a autonomia dos alunos na resolução de problemas matemáticos contextualizados, contribuindo para mitigar as dificuldades identificadas no ensino de matemática e oferecendo uma ferramenta alinhada às diretrizes da BNCC.

\textbf{Palavras-chave:} Pensamento Computacional. Gamificação. Engajamento. Resolução de Problemas. ENEM. Matemática.